% GENERATED FILE. Edit canonical lessons, then run npm run generate:books.
% canonical-source-hash: fnv1a64:4e31e710ff245b72
% canonical-lessons: RU-C465-liven, RU-C465-morosit, RU-C465-promokat, RU-C465-luzha, RU-C465-lit

\chapter{A downpour}
\label{ch:ru-a-downpour}
\hlchaptermodality{465}

\begin{chapteropening}
\textbf{By the end of this chapter:} \emph{I can talk about rain: a downpour, drizzle, getting soaked, puddles and rain pouring down.}

Complete the last lesson of chapter 465: I can talk about rain: a downpour, drizzle, getting soaked, puddles and rain pouring down.
\end{chapteropening}

\section[\texorpdfstring{\ru{ливень} (líven)}{ливень (líven)}]{\ru{ливень} (líven) — a downpour, a heavy shower}

\label{lesson:RU-C465-liven}



\begin{quote}
Before the new one: say the Russian for overcast, then the Russian for to clear up.
\end{quote}

\subsection*{You'll want to know: \ru{ливень}}
\textbf{\ru{ливень}} (\emph{líven}) — \textquotedblleft{}a downpour, a heavy shower\textquotedblright{}.

\begin{quote}
\ru{После} \ru{обеда} \ru{начался} \ru{сильный} \ru{ливень}.
\end{quote}

\begin{quote}
\emph{Pósle obéda nachalsyá sílnyy líven.}
\end{quote}

\begin{quote}
After lunch a heavy downpour began.
\end{quote}

Masculine; the \ru{е} drops out: \emph{\ru{ливня}}.

\begin{cousinweb}
\emph{\ru{ливень}} is built on \emph{\ru{лить}} 'to pour' with \emph{-\ru{ень}}: rain that pours. EtymDB takes \emph{\ru{лить}} from Proto-Slavic \emph{liti} 'to pour'.
\end{cousinweb}

\begin{grammarlens}[title={What you've built}]
One more word for telling, explaining and describing at B1.
\end{grammarlens}

\subsection*{Your turn}
\begin{itemize}
\raggedright
  \item \emph{Say it:} \emph{líven}
  \item \emph{Say it:} \emph{líven}, once more
  \item \emph{Say it:} a sentence of your own with \emph{líven}
  \item \emph{Recall:} say the Russian for overcast, then the Russian for to clear up, then say \emph{líven} again
\end{itemize}

\begin{tcolorbox}[breakable,colback=teal!4,colframe=teal!35!black,title={Before you move on}]
What does \ru{ливень} mean? (\textquotedblleft{}A downpour, a heavy shower\textquotedblright{}.) Say it once more.
\end{tcolorbox}
\section[\texorpdfstring{\ru{моросить} (morosít)}{моросить (morosít)}]{\ru{моросить} (morosít) — to drizzle}

\label{lesson:RU-C465-morosit}



\begin{quote}
Before the new one: say the Russian for to clear up, then the Russian for a downpour.
\end{quote}

\subsection*{You'll want to know: \ru{моросить}}
\textbf{\ru{моросить}} (\emph{morosít}) — \textquotedblleft{}to drizzle\textquotedblright{}.

\begin{quote}
\ru{С} \ru{утра} \ru{моросит} \ru{мелкий} \ru{дождь}.
\end{quote}

\begin{quote}
\emph{S utrá morosít mélkiy dozhd.}
\end{quote}

\begin{quote}
A fine rain has been drizzling since morning.
\end{quote}

Imperfective; mostly impersonal: \emph{\ru{моросит}} — it is drizzling.

\begin{cousinweb}
\emph{\ru{моросить}} is built on \emph{\ru{морось}} 'drizzle, fine rain'. EtymDB lists neither word, so its deeper origin is left open.
\end{cousinweb}

\begin{grammarlens}[title={What you've built}]
One more word for telling, explaining and describing at B1.
\end{grammarlens}

\subsection*{Your turn}
\begin{itemize}
\raggedright
  \item \emph{Say it:} \emph{morosít}
  \item \emph{Say it:} \emph{morosít}, once more
  \item \emph{Say it:} a sentence of your own with \emph{morosít}
  \item \emph{Recall:} say the Russian for to clear up, then the Russian for a downpour, then say \emph{morosít} again
\end{itemize}

\begin{tcolorbox}[breakable,colback=teal!4,colframe=teal!35!black,title={Before you move on}]
What does \ru{моросить} mean? (\textquotedblleft{}To drizzle\textquotedblright{}.) Say it once more.
\end{tcolorbox}
\section[\texorpdfstring{\ru{промокать} (promokát)}{промокать (promokát)}]{\ru{промокать} (promokát) — to get soaked, to get wet through}

\label{lesson:RU-C465-promokat}



\begin{quote}
Before the new one: say the Russian for a downpour, then the Russian for to drizzle.
\end{quote}

\subsection*{You'll want to know: \ru{промокать}}
\textbf{\ru{промокать}} (\emph{promokát}) — \textquotedblleft{}to get soaked, to get wet through\textquotedblright{}.

\begin{quote}
\ru{Без} \ru{зонта} \ru{мы} \ru{всегда} \ru{промокаем}.
\end{quote}

\begin{quote}
\emph{Bez zónta my vsegdá promokáyem.}
\end{quote}

\begin{quote}
Without an umbrella we always get soaked.
\end{quote}

Imperfective; perfective \emph{\ru{промокнуть}} (past \emph{\ru{промок}}). \emph{\ru{Промокнуть} \ru{до} \ru{нитки}} — to get soaked to the skin.

\begin{cousinweb}
\emph{\ru{промокать}} is built from \emph{\ru{про}-} 'through' and \emph{-\ru{мокать}}, on the root of \emph{\ru{мокрый}} 'wet', which is inherited from Proto-Slavic \emph{mokrŭ} (EtymDB).
\end{cousinweb}

\begin{grammarlens}[title={What you've built}]
One more word for telling, explaining and describing at B1.
\end{grammarlens}

\subsection*{Your turn}
\begin{itemize}
\raggedright
  \item \emph{Say it:} \emph{promokát}
  \item \emph{Say it:} \emph{promokát}, once more
  \item \emph{Say it:} a sentence of your own with \emph{promokát}
  \item \emph{Recall:} say the Russian for a downpour, then the Russian for to drizzle, then say \emph{promokát} again
\end{itemize}

\begin{tcolorbox}[breakable,colback=teal!4,colframe=teal!35!black,title={Before you move on}]
What does \ru{промокать} mean? (\textquotedblleft{}To get soaked, to get wet through\textquotedblright{}.) Say it once more.
\end{tcolorbox}
\section[\texorpdfstring{\ru{лужа} (lúzha)}{лужа (lúzha)}]{\ru{лужа} (lúzha) — a puddle}

\label{lesson:RU-C465-luzha}



\begin{quote}
Before the new one: say the Russian for to drizzle, then the Russian for to get soaked.
\end{quote}

\subsection*{You'll want to know: \ru{лужа}}
\textbf{\ru{лужа}} (\emph{lúzha}) — \textquotedblleft{}a puddle\textquotedblright{}.

\begin{quote}
\ru{Ребёнок} \ru{прыгал} \ru{по} \ru{лужам}.
\end{quote}

\begin{quote}
\emph{Rebyónok prýgal po lúzham.}
\end{quote}

\begin{quote}
The child was jumping in the puddles.
\end{quote}

Feminine.

\begin{cousinweb}
\emph{\ru{лужа}} is inherited from Proto-Slavic \emph{luža} 'a pool' (EtymDB).
\end{cousinweb}

\begin{grammarlens}[title={What you've built}]
One more word for telling, explaining and describing at B1.
\end{grammarlens}

\subsection*{Your turn}
\begin{itemize}
\raggedright
  \item \emph{Say it:} \emph{lúzha}
  \item \emph{Say it:} \emph{lúzha}, once more
  \item \emph{Say it:} a sentence of your own with \emph{lúzha}
  \item \emph{Recall:} say the Russian for to drizzle, then the Russian for to get soaked, then say \emph{lúzha} again
\end{itemize}

\begin{tcolorbox}[breakable,colback=teal!4,colframe=teal!35!black,title={Before you move on}]
What does \ru{лужа} mean? (\textquotedblleft{}A puddle\textquotedblright{}.) Say it once more.
\end{tcolorbox}
\section[\texorpdfstring{\ru{лить} (lit)}{лить (lit)}]{\ru{лить} (lit) — to pour}

\label{lesson:RU-C465-lit}



\begin{quote}
Before the new one: say the Russian for to get soaked, then the Russian for a puddle.
\end{quote}

\subsection*{You'll want to know: \ru{лить}}
\textbf{\ru{лить}} (\emph{lit}) — \textquotedblleft{}to pour\textquotedblright{}.

\begin{quote}
\ru{Дождь} \ru{льёт} \ru{с} \ru{утра}.
\end{quote}

\begin{quote}
\emph{Dozhd lyot s utrá.}
\end{quote}

\begin{quote}
It has been pouring since the morning.
\end{quote}

Imperfective. Present \emph{\ru{лью}, \ru{льёшь}, \ru{льют}}; past \emph{\ru{лил}, \ru{лила}}. \emph{\ru{Льёт} \ru{как} \ru{из} \ru{ведра}} — it is pouring with rain.

\begin{cousinweb}
\emph{\ru{лить}} is inherited from Proto-Slavic \emph{liti} 'to pour' (EtymDB).
\end{cousinweb}

\begin{grammarlens}[title={What you've built}]
One more word for telling, explaining and describing at B1.
\end{grammarlens}

\subsection*{Your turn}
\begin{itemize}
\raggedright
  \item \emph{Say it:} \emph{lit}
  \item \emph{Say it:} \emph{lit}, once more
  \item \emph{Say it:} a sentence of your own with \emph{lit}
  \item \emph{Recall:} say the Russian for to get soaked, then the Russian for a puddle, then say \emph{lit} again
\end{itemize}

\begin{tcolorbox}[breakable,colback=teal!4,colframe=teal!35!black,title={Before you move on}]
What does \ru{лить} mean? (\textquotedblleft{}To pour\textquotedblright{}.) Say it once more.
\end{tcolorbox}
